\section{Algorithm: Conservative Constrained Thompson Sampling}
\label{sec:algorithm}

We now develop \emph{Conservative Constrained Thompson Sampling} (CC-TS), a
four-layer policy that satisfies \cref{eq:bayesian-decision} while preserving
the conservative property under \cref{assumption:order}. The four layers are:

\begin{description}
  \item[Layer 1: Beta-Bernoulli posteriors.] Per-cell, per-tier Bayesian
        posteriors with offline-prior pseudocounts and online updates.
  \item[Layer 2: Conservative posterior-rejection gate.] One-sided
        credibility test that filters downgrades whose conformal lower bound
        on quality preservation falls below tolerance.
  \item[Layer 3: Asymmetric-loss decision rule.] Bayes-optimal action under
        the asymmetric loss matrix \cref{eq:loss}, with multipliers
        partitioned by tolerance class.
  \item[Layer 4: Uncertainty-triggered eval.] Eval suite is fired when
        posterior credibility is too wide for any decision to be confident.
\end{description}

The remainder of this section develops each layer, gives a regret bound for
the unconstrained limit, and proves the conservative-property guarantee.

\subsection{Layer 1: Beta-Bernoulli posteriors}
\label{sec:layer1}

For each cell $(\agent, \shape)$ and tier $\tier$, we maintain a
$\mathrm{Beta}(\alpha_{\agent,\shape,\tier}, \beta_{\agent,\shape,\tier})$
posterior on $\theta_{\agent,\shape,\tier}$. The Bernoulli likelihood
\cref{eq:bernoulli} is conjugate, so updates are closed-form: on observing
$q \in \{0,1\}$,
\[
  \alpha \leftarrow \alpha + q,
  \qquad
  \beta \leftarrow \beta + (1 - q).
\]

\paragraph{Cold-start prior.}
We seed each cell from the offline empirical landscape (the
agent-by-shape-by-tier matrix produced in \cite{blackrim_moa1b}). Each cell
in the landscape is annotated with a confidence percentage $c \in [0, 100]$
reflecting the strength of evidence underlying the recommendation. We map
this to Beta pseudocounts via
\begin{equation}
  \alpha_{\mathrm{prior}} = \max\!\big(1, \mathrm{round}(c / 5)\big),
  \quad
  \beta_{\mathrm{prior}} = \max\!\big(1, \mathrm{round}((100 - c) / 5)\big).
  \label{eq:prior-from-confidence}
\end{equation}
A 95\%-confidence cell maps to $\mathrm{Beta}(19, 1)$ (very informative); a
65\%-confidence cell maps to $\mathrm{Beta}(13, 7)$ (effectively
uninformative); a never-evidenced cell maps to $\mathrm{Beta}(1, 1)$, the
uniform prior.

\paragraph{Hierarchical partial pooling.}
Many cells will be sparse for weeks. We borrow strength across siblings
within an agent (different shapes, same agent) and across agents within a
tolerance class. The simplest closed-form approximation is weighted-average
pseudocount pooling:
\begin{equation}
  \tilde{\alpha} = \alpha + \lambda \sum_{\agent', \shape' \in \mathrm{sib}}
    w_{\agent',\shape'} \, \alpha_{\agent',\shape',\tier},
  \quad
  \tilde{\beta} \text{ analogously},
  \label{eq:pooling}
\end{equation}
with weights $w$ inversely proportional to sibling-posterior standard
error and a global cap $\lambda \leq 0.5$ to prevent pooling from dominating
own-cell evidence. Full hierarchical Beta-Beta inference is a
straightforward MOA-9 design choice (Stan \cite{carpenter2017stan} or PyMC
implementations exist); for the production loop the closed-form
approximation suffices and is provably no worse asymptotically.

\subsection{Layer 2: Conservative posterior-rejection gate}
\label{sec:layer2}

Given the pooled posterior, the conservative gate determines whether each
candidate tier is admitted to the feasible set
$\tierset_{\posterior, \tolerance, \alpha}$.

\paragraph{Conformal lower bound.}
We compute the conformal lower bound on quality preservation
\cite{angelopoulos2021gentle,vovk2005algorithmic}. A held-out calibration
set of $(predicted, observed)$ quality pairs is maintained per cell on a
rolling window. Conformal prediction emits a distribution-free interval
$[\widehat{q}_{\mathrm{lo}}, \widehat{q}_{\mathrm{hi}}]$ such that
\(
  \Pr(\theta \in [\widehat{q}_{\mathrm{lo}}, \widehat{q}_{\mathrm{hi}}])
    \geq 1 - \alpha,
\)
regardless of the underlying distribution. The lower bound
$\widehat{q}_{\mathrm{lo}}(\agent, \shape, \tier)$ is then plugged into the
gate:

\begin{definition}[Conservative gate]
\label{def:gate}
For a candidate tier $\tier \prec \tier^*$ at cell $(\agent, \shape)$ with
tolerance $\tolerance$, the conservative gate \emph{admits} $\tier$ if and
only if
\begin{equation}
  \widehat{q}_{\mathrm{lo}}(\agent, \shape, \tier)
    \;\geq\;
  \mu^*(\agent, \shape) - \tolerance,
  \label{eq:gate}
\end{equation}
where $\mu^*$ is the posterior-mean quality at the baseline tier.
\end{definition}

The gate has four operational regimes, controlled by tolerance class:
$\tolerance = 0$ (Critical) admits no downgrade ever, regardless of evidence
strength; $\tolerance < 0.02$ (Strict) admits a downgrade only when the
conformal LCB exceeds baseline-mean by 2~percentage points; and so on
through Lenient.

\subsection{Layer 3: Asymmetric-loss decision rule}
\label{sec:layer3}

Among the admitted tiers, we pick the action minimizing expected
asymmetric loss. The loss matrix, indexed by action
$a \in \{\mathrm{down}, \mathrm{stay}, \mathrm{up}\}$ and outcome $\omega
\in \{\mathrm{preserved}, \mathrm{lost}\}$, is shown in
\cref{tab:loss-matrix}.

\begin{table}[h]
\centering
\caption{Asymmetric loss matrix. The Critical multiplier $M[\text{Critical}] = \infty$
collapses to a hard rejection: critical-path cells never enter the
exploration set.}
\label{tab:loss-matrix}
\begin{tabular}{lll}
\toprule
\textbf{Action $\times$ outcome} & \textbf{Loss} & \textbf{Notes} \\
\midrule
$L(\mathrm{down}, \mathrm{preserved})$ & $-\Delta\cost(\tier^*, \tier_{\mathrm{lower}})$ & savings \\
$L(\mathrm{down}, \mathrm{lost})$
  & $\Delta\cost_{\mathrm{down}} \cdot N_{\mathrm{dep}} \cdot M[\tolerance]$
  & cascade \\
$L(\mathrm{stay}, \cdot)$ & $0$ & baseline \\
$L(\mathrm{up}, \mathrm{preserved})$ & $+\Delta\cost(\tier_{\mathrm{higher}}, \tier^*)$ & overpay \\
$L(\mathrm{up}, \mathrm{lost})$ & $+\Delta\cost(\tier_{\mathrm{higher}}, \tier^*)$ & overpay \\
\bottomrule
\end{tabular}
\end{table}

The multipliers $M[\cdot]$ are
\(
  M[\text{Critical}] = \infty,
\)
\(
  M[\text{Strict}] = 20,
\)
\(
  M[\text{Moderate}] = 5,
\)
\(
  M[\text{Lenient}] = 1.
\)
A Critical cell with $M = \infty$ degenerates the decision rule to
``never downgrade,'' which is the operational behaviour we want for ADRs,
threat models, and incident-response orchestration. Lenient cells
($M = 1$) treat downstream cascade as bounded by the saved API spend, which
is the correct calculus for, e.g., a research lookup that the human
operator gates anyway.

\subsection{Layer 4: Uncertainty-triggered eval}
\label{sec:layer4}

When the conformal interval at the chosen tier is wider than a threshold
$\theta_{\mathrm{eval}}$ (we use $\theta_{\mathrm{eval}} = 0.10$ in
production), the advisor schedules a deterministic eval-suite run on the
cell. Eval results are weighted higher in the likelihood than production
quality observations because their ground truth is known; the conformal
calibration set is updated with the eval $(predicted, observed)$ pairs.

This is the loop that closes the regulator: cells with thin evidence ---
notably the open-questions cells \texttt{OQ-1}--\texttt{OQ-5} from
\cite{blackrim_moa1b} --- accumulate eval probes proportional to their
posterior width, and the advisor's policy converges on those cells in time
proportional to eval cadence rather than production traffic.

\subsection{The CC-TS algorithm}

\Cref{alg:ccts} states CC-TS in full.

\begin{algorithm}[h]
\KwIn{Cell $(\agent, \shape)$, tolerance $\tolerance$, baseline tier $\tier^*$,
      pooled posterior $\widetilde{\posterior}$,
      conformal calibration sets $\mathcal{C}_{\agent,\shape,\tier}$,
      loss matrix $L$, eval threshold $\theta_{\mathrm{eval}}$,
      confidence level $\alpha = 0.05$.}
\KwOut{Recommended tier $\tier^\dagger$, rationale $r$.}

\BlankLine
$\mathrm{cands} \leftarrow \{\tier^*\}$ \tcp*{baseline always admitted}
\For{$\tier \in \tierset \setminus \{\tier^*\}$}{
  $\widehat{q}_{\mathrm{lo}} \leftarrow
    \mathrm{ConformalLCB}(\widetilde{\posterior}_{\tier},
    \mathcal{C}_{\agent,\shape,\tier}, \alpha)$\;
  \If{$\tier \prec \tier^*$ \KwAnd $\widehat{q}_{\mathrm{lo}}
       \geq \mu^*(\agent, \shape) - \tolerance$}{
    $\mathrm{cands} \leftarrow \mathrm{cands} \cup \{\tier\}$
      \tcp*{Layer 2: conservative gate}
  }
  \ElseIf{$\tier \succ \tier^*$}{
    $\mathrm{cands} \leftarrow \mathrm{cands} \cup \{\tier\}$
      \tcp*{upgrades unconstrained}
  }
}

$\tier^\dagger \leftarrow
  \arg\min_{\tier \in \mathrm{cands}} \;
    \mathbb{E}_{\theta \sim \widetilde{\posterior}_{\tier}}
      \big[ L(\tier, \theta; \tier^*, \mu^*) \big]$
  \tcp*{Layer 3: asymmetric-loss min}

\If{$\mathrm{CIWidth}(\widetilde{\posterior}_{\tier^\dagger}, \alpha)
     > \theta_{\mathrm{eval}}$}{
  $\mathrm{ScheduleEval}(\agent, \shape, \tier^\dagger)$
    \tcp*{Layer 4: uncertainty-triggered}
}

\Return{$\tier^\dagger,
  r = \{\mathrm{cands}, \widehat{q}_{\mathrm{lo}}, \mathrm{LossExp}\}$}\;

\caption{CC-TS recommend.}
\label{alg:ccts}
\end{algorithm}

\subsection{Theoretical guarantees}

We give two results: a regret bound under unconstrained Thompson Sampling
(when the conservative gate is vacuous because every tier passes), and a
conservative-property guarantee that holds always.

\begin{proposition}[Unconstrained regret]
\label{prop:regret}
In the unconstrained limit ($M = 0$, equivalently $\tolerance \geq 1$, every
candidate admitted), CC-TS reduces to standard Thompson Sampling per cell.
By Agrawal--Goyal \cite{agrawal2012analysis}, the per-cell expected regret
after $T$ rounds satisfies
\begin{equation}
  \mathbb{E}[R_T] \;\leq\; \mathcal{O}\!\left(
    \sum_{\tier \neq \tier^\star_{\agent,\shape}}
      \frac{\log T}{\Delta_\tier}
  \right),
  \label{eq:regret}
\end{equation}
with $\Delta_\tier = \theta_{\agent,\shape,\tier^\star} - \theta_{\agent,\shape,\tier}$
the per-arm gap.
\end{proposition}

The proof is the standard Agrawal--Goyal decomposition; we omit it.

\begin{theorem}[Conservative property]
\label{thm:conservative}
Suppose conformal calibration is exchangeable with the production stream
\cite{vovk2005algorithmic}. Then for any cell $(\agent, \shape)$ and
tolerance $\tolerance$, CC-TS recommends a downgrade
$\tier \prec \tier^*$ only if
\begin{equation}
  \Pr\!\big[
    \theta_{\agent,\shape,\tier} \geq \theta_{\agent,\shape,\tier^*} - \tolerance
  \big]
  \;\geq\;
  1 - \alpha.
  \label{eq:conservative}
\end{equation}
\end{theorem}

\begin{proof}[Proof sketch.]
The downgrade $\tier \prec \tier^*$ is admitted to $\mathrm{cands}$ in
\cref{alg:ccts} only if the conformal lower bound exceeds
$\mu^*(\agent, \shape) - \tolerance$. By the marginal coverage guarantee of
conformal prediction \cite{angelopoulos2021gentle}, the event
$\theta_{\agent,\shape,\tier} \geq \widehat{q}_{\mathrm{lo}}$ has probability
at least $1 - \alpha$ under the calibration distribution. Combining the two
yields \cref{eq:conservative}. The asymmetric-loss step (Layer 3) can only
remove tiers from the candidate set, never add them, so the inequality is
preserved by the final selection.
\end{proof}

\Cref{thm:conservative} is the load-bearing property of the system: every
operationally-relevant downgrade is credible at the chosen confidence
level, and the credibility is wired to the same conformal machinery whose
coverage is well-understood \cite{angelopoulos2021gentle}.

\subsection{Implementation notes}

The reference implementation lives in \texttt{internal/dispatch/}
\cite{german2026blackrim}. Three implementation properties are worth
flagging because they materially affect the deployment story:

\begin{enumerate}[leftmargin=*]
  \item \textbf{Pure-function recommend.} The function
        \texttt{Recommend(state, key) -> (tier, reasons)} has no I/O and no
        global mutable state, which makes it trivial to test exhaustively.
  \item \textbf{Determinism for tests.} The Beta sampler is seedable from a
        \texttt{Recommend} argument, so unit tests can verify decision shapes
        across pseudo-random draws.
  \item \textbf{Reasons object.} Every recommendation returns a structured
        rationale: candidate tiers, conformal LCBs, expected losses, and
        whether eval was triggered. This is the audit surface and the
        operator interface in one structure.
\end{enumerate}

The full pseudocode in
\cite{blackrim_moa1c} is faithful to the implementation; we elide it here
for space.
